\subsection*{Two-coordinate A.11 interval-net certification and a strict rational margin (22 October 2026)}
\label{subsec:a11-two-coordinate-net-and-margin}
This checkpoint supplements the previously derived exact terminal-gauge
residual lift. The full corrected A.11 hull \(\mathcal H\) contains
both azimuthal supports \emph{and} the intervening gap and depends on
the pulse coordinate \(y\) and angular coordinate \(\eta\).
An angular grid alone at fixed \(y\) is insufficient to establish a
uniform estimate on \(\mathcal H\).

Fix the \emph{same} authentic corrected A.11/A.19 profile and define
\[
G_\pm(y,\eta)=\lambda E(y,\eta)\pm 2N_s(y,\eta).
\]
When \(E>0\), the exact desired source stress criterion is equivalent to
\(G_+\ge0\) and \(G_-\ge0\). There is no need to assume that
the physical source integrals are polynomial in \(\eta\).

\begin{proposition}[Two-coordinate interval-net gate]
Let \(V\subseteq\mathcal H\) be finite. Suppose every
\(z=(y,\eta)\in\mathcal H\) has a node \(v=(y_i,\eta_j)\in V\) with
\[
|y-y_i|\le r_y,\qquad|\eta-\eta_j|\le r_\eta.
\]
For each \(f\in\{E,G_+,G_-\}\), suppose independently certified
nonnegative constants \(L_{f,y},L_{f,\eta}\) satisfy
\[
|f(z)-f(v)|\le L_{f,y}|y-y_i|+
                    L_{f,\eta}|\eta-\eta_j|.
\]
If for a common \(\delta>0\) every node satisfies
\[
E(v)\ge\delta+L_{E,y}r_y+L_{E,\eta}r_\eta,\qquad
G_\pm(v)\ge L_{\pm,y}r_y+L_{\pm,\eta}r_\eta,
\]
then \(E(z)\ge\delta>0\) and \(G_\pm(z)\ge0\)
for every \(z\in\mathcal H\). In particular
\(\sup_{\mathcal H}|N_s|/E\le\lambda/2\).
\end{proposition}
\begin{proof}
For each point choose a covering node and subtract the
two coordinatewise Lipschitz-remainder envelopes. The two signed
inequalities imply \(2|N_s|\le\lambda E\), and \(E>0\)
permits division.
\end{proof}
The \(\eta\)-derivative inputs must be \emph{total} derivatives along the
same genuine A.11/A.19 source root, while \(y\)-derivative inputs must
cover support boundaries, the gap, and applicable stage transitions.
A grid scan without authenticated derivative/remainder bounds is
\emph{not} a uniform certificate.

\begin{proposition}[Explicit release margin for the residual lift]
Assume the earlier exact same-profile residual inequality
\(\sup_{\mathcal H}|N_s|/E\le b_0+\mathcal R\),
\(E>0\), and \(\lambda>0\), where
\[
\mathcal R=\alpha[(2+4h)\|\mathcal F\|_\infty+
\|\mathcal F_\eta\|_\infty]+
\beta[(1+2h)\|\mathcal G\|_\infty+
\tfrac12\|\mathcal G_\eta\|_\infty]
\]
with the same definitions and authenticated premises as in the
preceding source-gauge residual section. If
\[
b_0\le\frac{\lambda}{500},\qquad
\mathcal R\le\frac{4979\lambda}{10000},
\]
then the source stress ratio has the strict spare margin
\[
\boxed{\quad
\sup_{\mathcal H}\frac{|N_s|}{E}
\le\frac{\lambda}{2}-\frac{\lambda}{10000}
<\frac{\lambda}{2}.\quad}
\]
\end{proposition}
\begin{proof}
Exactly \(1/500+4979/10000=4999/10000
=1/2-1/10000\). Positivity of \(\lambda\) gives strictness.
\end{proof}

\paragraph{Source authentication and limits.}
Both propositions are complete conditional scalar arguments.
They do not supply the true \((y,\eta)\) intervals or root derivatives.
The companion exact-rational verifier replays
10,000 scalar budgets and 10,000 two-coordinate affine controls and
rejects three negative controls. The Lean source
\texttt{A11RectangularNetScalar.lean} formalizes the corresponding
finite implications subject to independent CI success; it does not
prove the actual PDE-specific source bounds. The authentic joint
parameter tuple and A.11/A.19 roots, full stress cones, Appendix B
contraction, unforced startup hierarchy, physical zero-force
Navier--Stokes residual and the proposed Theorem~1.1 all remain
OPEN/HOLD. The original OpenAI construction has smooth forcing.
